\section{Discussion}
The reviewed studies are concentrated around AQI classification, risk alerting, and pollutant prediction, and 13 of the 19 studies report real-time or near-real-time operation. At the same time, cross-study comparison remains limited because the corpus varies in input pollutants, fuzzy-system family, deployment architecture, and evaluation metrics. These results are in line with AQI antecedents that describe heterogeneity in index construction and interpretation \cite{ART01,ART02,PROB08}, and with applied fuzzy-monitoring studies that address related problems through different platforms, contexts, and objectives \cite{ART04,ART05,ART06,ART07,ART08}. The discussion therefore focuses on the main patterns observed in the corpus and on their implications for comparability, transferability, and future platform design.

\subsection{Insufficient External Generalization}
One of the clearest findings of the review is that the good performance reported by studies does not automatically translate into solid evidence of external generalization. In the corpus, several works achieve favorable results within their own experimental environments, but still maintain limited validation when the city, the sensor network, or the operational window changes \cite{A028,A084,A087,A144}. This pattern can be observed, for example, in \cite{A028}, where AQI is predicted from public data but without operational deployment or verification in more than one city; in \cite{A144}, which acknowledges limited deployment scope and the need for validation at multiple sites; and in \cite{A084}, where adjustment sensitivity varies by city and requires additional validation.

The applied antecedents point in the same direction: \cite{ART04} uses an open dataset with a classification focus and proposes expanding test locations as future work; \cite{ART05} reports real-time monitoring, but without validation supported by analyzed and comparable datasets; \cite{ART06} and \cite{ART07} implement IoT monitoring with fuzzy logic and real-time operation, although without sufficient traceability of external data for reproducible contrast; and \cite{ART08} advances further in openness by publishing dataset and code, in addition to cross-validation, but still within a local deployment context. Both the corpus and these antecedents therefore show that the functional viability of the approach is better documented than its robust transferability across different contexts.

Our corpus adds another nuance: this weakness does not depend only on sample size or type, but also on the combination of data source, evaluation form, and level of methodological reporting. In the corpus, studies based on direct sensors predominate (\texttt{13/19}) over hybrid sources (\texttt{5/19}) and public-only sources (\texttt{1/19}), while real-time or near-real-time operation is frequent (\texttt{13/19}). However, reporting gaps persist (\texttt{4/19} without clearly declared limitations), together with specific cases in which input information was not explicitly reported in the extraction (A028 and A084 with \texttt{D1\_Reported\_input\_values=NR}) and a heterogeneous evaluation landscape (\texttt{Accuracy=8}, \texttt{RMSE=7}, \texttt{MAE=5}, \texttt{MAPE=5}, \texttt{R2=3}) that hinders direct comparisons. Methodologically, this indicates that generalization should not be understood only as replication in more than one site, but as the result of multi-context validation accompanied by traceable data and comparable protocols.

\subsection{Limited Coverage of Variables and Pollutants}
Another consistent pattern is the concentration of modeling around a relatively limited set of pollutants and support variables, even though the air-quality phenomenon is chemically broader and more complex. In D1, the most recurrent pollutants are PM$_{2.5}$ and CO (eight studies each), followed by CO$_2$ and PM$_{10}$, while temperature and humidity predominate among environmental variables. The studies also distribute these inputs differently according to output objective: PM$_{2.5}$ is recurrent in forecasting and calibration settings \cite{A084,A135,A144,A157}, whereas CO, CO$_2$, and PM$_{10}$ appear more often in AQI/IAQ classification, index estimation, and multi-gas configurations \cite{A009,A028,A136,A149,A167,A168,A178}. The concentration therefore describes what is measured most often, but also which output classes receive greater methodological attention in the reviewed platforms.

The reviews by Priti and Kumar \cite{ART01} and Suman \cite{ART02}, as well as Murena's urban index development \cite{PROB08}, had already shown that the representation of air-quality status depends on which pollutants enter the index, how they are aggregated, and under what severity interpretation criterion. What now appears in this SLR is the same problem transferred to the plane of fuzzy platforms, where not only index models change, but also the subset of pollutants incorporated into monitoring and inference, conditioning system sensitivity to different risk profiles and favoring partial readings when only a few variables predominate.

Something similar occurs in the applied antecedents on fuzzy logic, where the scope of variables depends strongly on deployment context and system objective. The classification, environmental control, and IoT monitoring studies reviewed in \cite{ART04,ART05,ART06,ART07} work with rather contained sets of inputs that are functional to their use case, whereas \cite{ART08} incorporates the dimension of distributed operation and network architecture, but also does not broadly redefine pollutant coverage. What our corpus adds is that this reduction responds not only to context, but also to practical commitments related to instrumentation, data availability, and processing cost; for this reason, parsimonious configurations based on few pollutants and recurrent microclimatic covariates end up predominating.

The cross-cutting indicators point in the same direction: in \texttt{11/19} studies (\texttt{57.9\%}), at most two pollutants are modeled (\texttt{1}=7, \texttt{2}=4), whereas only \texttt{6/19} work with four or more. In addition, \texttt{5/19} do not report support environmental variables, and direct acquisition predominates (\texttt{13/19}). The most recurrent combination links particulate matter, CO/CO$_2$, temperature, and humidity, whereas O$_3$, NOx, NH$_3$, CH$_4$, smoke, and specialized particle counts appear in more specific settings. The corpus therefore leans toward compact input configurations, while leaving less coverage for other pollutant profiles and operating contexts.

\subsection{Integration and Deployment under Operational Conditions}
Another relevant result is that the corpus does not consist only of offline models: \texttt{13/19} studies report real-time or near-real-time operation, \texttt{9} include dashboards, \texttt{10} describe continuous-loop behavior, and \texttt{5} incorporate alarm modules. Here, \emph{operation} refers to the combined functioning of data acquisition, processing or inference, visualization, and, when applicable, alerting or actuation under actual use conditions. Even so, these components are not always reported as a fully connected end-to-end flow. This can be seen in \cite{A028}, which reports strong model fit with public data but no operational deployment; in \cite{A135}, which documents PM$_{2.5}$ forecasting performance without full integration of the forecasting module into the online portal; and in \cite{A144}, which reports predictive performance with limited deployment scope and the need for multi-site validation.

The applied antecedents describe comparable scenarios. \cite{ART05} and \cite{ART07} come closer to real-time functioning by combining sensing, inference, and response over the monitored environment; \cite{ART06} also articulates IoT monitoring and index interpretation; and \cite{ART08} broadens the panorama by incorporating distributed architecture and network-efficiency dimensions for smart-city scenarios. In both fronts, antecedents and corpus, systems already exist that connect sensing, inference, and operational output; what continues to vary from one study to another is the degree of end-to-end integration and the level of detail with which that integration is reported.

Our corpus also adds a documentation problem: the difficulty lies not only in deploying components, but in describing their coupling in a comparable way. Within the corpus, \cite{A114} stands out as a positive reference by integrating monitoring, visualization, and actuation with reported inference latency. By contrast, in other studies, the presence of dashboards, continuous loops, or online inference is not sufficient by itself to reconstruct how acquisition, processing, storage, publication, and response are connected. For cross-study comparison, speaking about operation should therefore imply something more than declaring real time or dashboard presence: it should allow reconstruction of how the system acquires, processes, publishes, and sustains its functioning in context.

\subsection{Computational Cost and Network Efficiency}
The analysis of D3 and evaluation reveals another imbalance: predictive and classificatory performance are reported more often than computational cost or network efficiency. This can be seen in \cite{A135}, which reports forecasting performance together with high training time and incomplete operational integration; in \cite{A138}, which compares optimized variants in terms of accuracy and computation time; and in \cite{A178}, which leaves network efficiency and user-interaction performance as pending work. These cases show that cost and efficiency are considered in some studies, but much less consistently than predictive or classification performance.

The applied antecedents point in the same direction. Among them, \cite{ART08} comes closest to an explicit treatment of network efficiency, whereas \cite{ART05,ART06,ART07} focus more on monitoring, classification, and functional response. The corpus follows this pattern: although \texttt{13/19} studies report real-time or near-real-time operation, only \texttt{4/19} include metrics from the \textit{System} family. In addition, the reported efficiency indicators are dispersed and rarely repeated (\texttt{Training\_Time}, \texttt{Inference\_Time}, \texttt{Temporal\_Delay}, \texttt{Throughput}, \texttt{Reliability}, \texttt{Stability\_AAS}), which makes direct comparison difficult. Under these conditions, high fit or classification accuracy does not by itself indicate how demanding the platform is in sustained computational or communication terms.

\subsection{Hardware, Sensor, and Calibration Constraints}
The corpus also shows that platform performance is tied to instrumentation and sensing conditions. Relatively standard instrumentation chains predominate, with \texttt{PM\_Sensor}=10, \texttt{TempHumidity\_Sensor}=9, \texttt{CO\_Sensor}=8, \texttt{MCU\_Arduino}=7, \texttt{CO2\_Sensor}=5, and \texttt{VOC\_Sensor}=4, together with a predominance of direct acquisition (\texttt{13/19}). This configuration indicates recurring choices in sensors and hardware, while also leaving calibration, drift, precision, and sensing stability as relevant factors for comparing studies.

In the corpus, this dependence appears explicitly in \cite{A001}, where operational limits associated with platform and sensing are reported, including payload, autonomy, and wind sensitivity in UAVs, together with sensor nonlinearity, sensor precision, and communication range; and in \cite{A087}, where the system improves predictive performance for filter management in an industrial environment, but still acknowledges sensor drift and aging, the need for more robust calibration, and the absence of a closed criterion for the optimal number of sensors. Technical performance may therefore be favorable and still depend on instrumental assumptions that reduce transferability and stability outside the test environment.

The same dependence between system performance and sensing/infrastructure conditions is also evident in the presented antecedents. \cite{ART05}, \cite{ART06}, and \cite{ART07} describe functional real-time monitoring and classification systems supported by sensors and IoT platforms, whereas \cite{ART08} adds the network and distributed-operation dimension. What our review makes visible is that this aspect should not be treated as a peripheral hardware detail, because it ends up conditioning comparability across studies when calibration and maintenance procedures are not reported with sufficient homogeneity.

For this reason, comparison across fuzzy platforms should rest not only on output metrics or logical architecture, but also on how calibration, instrumental stability, aging, and sensing operating conditions are documented. Without that level of detail, system reproducibility becomes better established at the algorithmic level than at the physical acquisition level.

\subsection{Predictive-Horizon Degradation}
Performance stability when the temporal prediction horizon is extended appears as a problem much less resolved than classification or immediate estimation. In the corpus, this issue appears explicitly in \cite{A157}, where the fuzzy ARTMAP approach for PM$_{2.5}$ prediction with online training maintains acceptable global fit, but reports degradation at longer horizons because of error propagation. The point is relevant because, beyond the validity of the fuzzy approach itself, it questions temporal robustness when operational decision-making requires broader and more stable anticipation.

The review also shows that this issue can be examined only within a limited part of the corpus, because explicit PM$_{2.5}$ prediction appears in only \texttt{3/19} studies (\cite{A135,A144,A157}) and with different protocols (\texttt{TrainTest+CrossSite}, \texttt{CV+Holdout}, \texttt{Online}). By comparison, the corpus gives greater weight to AQI classification, index estimation, and near-real-time monitoring than to multi-horizon predictive evaluation.

The applied antecedents point in the same direction. Their main focus lies on classification, real-time monitoring, local control, or architectural deployment \cite{ART04,ART05,ART06,ART07,ART08}, rather than on error behavior across predictive horizons. Within the corpus, predictive tasks are also difficult to compare directly because the target pollutant, validation protocol, and temporal scale vary from one study to another.

Looking ahead to future studies, fuzzy prediction for air quality should be evaluated not only through aggregate error or global fit, but also through explicit temporal stability according to prediction horizon. Without this latter dimension, comparison among predictive models may hide the fact that the same system can be useful for immediate anticipation and, at the same time, unstable when the operational horizon is extended.

\subsection{Clinical Domain without Strong Causal Closure}
In studies with clinical or health-oriented focus, the review identifies a clear operational usefulness, but still with weak causal closure between air-quality indicators and health outcomes. In \cite{A142}, the fuzzy operating-room system classifies contamination-risk groups through interpretable rules and real-time alerts, but its own limitations acknowledge the absence of direct correlation with clinical infection incidence and the need to reconfigure rules as new variables are incorporated. In a convergent way, \cite{A149} improves FIAQI estimation through ARIMA-Kalman+FIS fusion over CO$_2$/TVOC, although it reports that long-term health evaluation by parameter remains open. In both cases, inference is useful for monitoring and prevention, but cause-effect traceability at the level of clinical event remains open.

The review also shows that this domain occupies a limited space within the corpus: the \texttt{Healthcare\_Facility} context appears in only \texttt{2} studies and the clinical unit \texttt{OR\_Risk\_Classification} in \texttt{1/19}. The comparative basis is therefore much smaller than in AQI classification, urban monitoring, or pollutant prediction, and the discussion of clinical applications has to be read within that narrower evidence base.

The applied antecedents reinforce this reading. In \cite{ART05,ART06,ART07}, a functional orientation toward monitoring, control, and interpretation of air status predominates, whereas \cite{ART08} concentrates on architecture and distributed operation; none of these antecedents proposes strong clinical causal validation. The same happens within the corpus: even when studies approach health-related environments, the relationship among environmental measurement, fuzzy risk classification, and clinical outcome still rests more on operational plausibility than on direct causal verification.

For this reason, health-oriented fuzzy platforms should differentiate more clearly between preventive utility, risk classification, and causal validation. Without this distinction, a system may be valuable as support for clinical monitoring and, at the same time, still not provide sufficient evidence to sustain strong inferences about health effects or medical outcomes.

\subsection{Review Limitations}
This review also presents limitations that should be made explicit. First, the documentary scope was restricted to Scopus, Web of Science, and IEEE Xplore, to English-language publications, and to the 2021--2026 period. Although this delimitation was coherent with the PICo framework and made it possible to work with a recent and manageable corpus, relevant studies indexed in other sources, gray literature, or works published outside this temporal window may have been left out.

Second, the synthesis depended on the level of detail reported by the primary studies. Although extraction was designed with page-based traceability and a controlled coding catalog, final comparability remained conditioned by heterogeneities that the review itself cannot eliminate, such as differences in output objectives, metrics, validation protocols, technical-detail level, and reporting completeness. Consequently, several findings should be read as consistent patterns within the analyzed corpus rather than as exhaustive generalizations about the whole field of air-quality monitoring with fuzzy logic.

\begin{adjustwidth}{-\extralength}{0cm}
\begin{longtblr}
[
 label={tab:cierre_vacios_implicaciones},
 caption={Analytical closure: recurrent gaps, methodological impact, and operational implications. Own elaboration.}]{
  colspec={X[2.2,l,m] X[0.8,c,m] X[0.8,c,m] X[0.9,c,m] X[2.6,l,m] X[2.8,l,m]},
  width=\fulllength,
  rowhead=1,
  row{1}={font=\bfseries}
}
\toprule
Recurrent gap/limit & k & n & \% & Impact on synthesis & Implication for the proposal \\
\midrule
Unreported data partition \texttt{NR} & 10 & 19 & 52.6 & Introduces uncertainty into performance comparison due to uncontrolled variation in the validation scheme. & Define and justify a fixed evaluation scheme (\texttt{train/val/test} or cross-validation), with full traceability of its application. \\
Unreported defuzzification & 8 & 19 & 42.1 & Hinders reproducibility of the inferential core and technical attribution of performance to fuzzy design. & Specify the defuzzification method and associated parameterization as a required pipeline component. \\
Missing comparator \texttt{NR} & 4 & 19 & 21.1 & Reduces the comparative validity of improvement claims against alternative approaches. & Include a minimum benchmark and a homogeneous criterion for quantitative comparison in experimentation. \\
Unreported improvement \texttt{NR} & 4 & 19 & 21.1 & Limits interpretation of incremental contribution by not making the magnitude of gain explicit. & Report absolute and relative improvement over the baseline under equivalent metrics and conditions. \\
Metric heterogeneity (without a single core) & 11 & 19 & 57.9 & Fragments cross-study reading and increases reporting heterogeneity across studies. & Establish a base metric set (\texttt{MSE/RMSE/MAE} and, when applicable, \texttt{Accuracy}) for mandatory reporting. \\
\bottomrule
\end{longtblr}
\end{adjustwidth}

Table~\ref{tab:cierre_vacios_implicaciones} helps condense where the obstacles that most affect the cross-reading of the corpus are concentrated. The highest weights fall on metric heterogeneity and the absence of reported data partition, followed by omission of defuzzification, three gaps that directly interfere with the possibility of comparing results, reconstructing design decisions, and attributing performance to concrete technical components. This pattern is consistent with what was discussed in the antecedents: AQI reviews had already warned about comparability problems in risk representation \cite{ART01,ART02}, and the applied antecedents showed functional solutions but reported with unequal levels of detail \cite{ART04,ART05,ART06,ART07,ART08}. What the corpus synthesis now shows is that these difficulties do not disappear when the focus shifts to real-time fuzzy platforms; rather, they reappear in the form of fragmented metrics, incomplete validation, and uneven technical description. From this point, the design implications presented below become easier to understand, because they are oriented toward reinforcing traceability, comparability, and generalization capacity.

\subsection{Design Implications Derived from D2 and D3}
The design implications derived from the review can be organized into three complementary planes: decisions about the fuzzy core, decisions about the operational architecture, and reproducibility criteria for a future proposal.

On a first plane, the evidence from D2 provides criteria for clearly delimiting the fuzzy component of the solution:
\begin{enumerate}
\item \textbf{Fuzzy role}: fuzzy inference appears primarily as the decision core rather than as auxiliary post-processing \cite{A001, A009, A149, A165}.
\item \textbf{Inference objective}: AQI classification, risk alerting, and pollutant prediction predominate, so output definition should remain aligned with the operational objective of use \cite{A009, A135, A165, A168}.
\item \textbf{Rule definition}: explicit rule bases and adaptive approaches are both observed, making it important to document when each strategy is adopted \cite{A132, A149, A144}.
\item \textbf{Membership functions}: triangular and trapezoidal linguistic functions predominate, with direct implications for semantic traceability \cite{A001, A138, A142}.
\item \textbf{Defuzzification}: the choice between \texttt{Centroid} and \texttt{Weighted average} appears linked to the FIS family and output type \cite{A001, A009, A123, A149}.
\item \textbf{Risk modeling}: AQI thresholds predominate, whereas customized scales require a more explicit justification \cite{A001, A114, A142, A178}.
\item \textbf{Interpretability}: transparency of rules, membership functions, and AQI mapping emerges as a recurrent auditability criterion \cite{A149, A165, A167}.
\end{enumerate}

On a second plane, the evidence from D3 points to an operational architecture with a common base and adaptable components. The most recurrent patterns in the corpus are concentrated in \texttt{IoT edge-cloud} or \texttt{Two-layer IoT} architectures, with real-time or near-real-time operation and with minimum documentation of connectivity, processing, and visualization. On that basis, hybrid or distributed strategy, instrumentation, and deployment type may vary, provided that each variant preserves sufficient traceability of output objective, metrics used, data partition, and latency or operating frequency.

On a third plane, these findings can be translated into reproducibility and comparability guidelines for a future proposal:
\begin{enumerate}
\item \textbf{Open-data traceability}: it would be advisable to record source, retrieval date/time, geographic coverage, pollutants used, and version of the extracted set.
\item \textbf{Near-real-time operation}: the pipeline should execute ingestion, preprocessing, inference, and publication with configurable frequency, reporting observed end-to-end latency conditioned by source update frequency.
\item \textbf{Auditable fuzzy core}: the rule base, membership functions, inference method, and defuzzification should be documented explicitly in order to ensure technical reproducibility.
\item \textbf{Comparable evaluation}: the proposal should report comparable unit, validation protocol, data partition, and a minimum set of metrics, avoiding contrasts between non-equivalent objectives.
\item \textbf{Multi-context external validation}: it would be advisable to verify performance in more than one city and in more than one temporal window, in order to assess stability outside the initial context.
\item \textbf{Data-quality assurance}: since physical sensors are not under control, it would be important to include systematic checks of completeness, consistency, outliers, and temporal stability in the input series.
\item \textbf{Operational efficiency}: it would also be advisable to report execution costs (processing/inference time and computational consumption), together with their impact on update frequency and scalability.
\end{enumerate}
